\subsubsection*{Solution}
\paragraph*{Step-by-Step Derivation}

The depolarized-GHZ fidelity convention and the model $F(n)=Fk^{n-1}$ match the formulation analyzed in the primary distributed-sensing study by Zang et al. \href{\detokenize{https://arxiv.org/abs/2409.17089}}{Quantum Advantage in Distributed Sensing with Noisy Quantum Networks}.

{\color{blue}

The AI agent found the work based on which (I suppose) the question has been formulated. In fact the formulation of the question is so close to the sentences taken from this work that I believe the AI agent did not even try to "think"/deviate from the original derivation. For this reason:
\begin{enumerate}
  \item The ambiguity of the assumption about the GHZ state fidelity has not been questioned and commented about. The author of the work specifically state it as a "phenomenological model" that is non-standard in the community (they even had to write an appendix to justify it).
  \item The presented derivations are as convoluted as those in the referenced work.
\end{enumerate}

}

\subparagraph*{1. Initial depolarized GHZ state}

Let

\[
N=nd,\qquad D=2^N,\qquad f=F(n)=Fk^{n-1}.
\]

Define

\[
|\mathrm{GHZ}_N^{+}\rangle = \frac{|0\rangle^{\otimes N}+|1\rangle^{\otimes N}}{\sqrt2}.
\]

A depolarized GHZ state whose actual GHZ fidelity is $f$ is

\[
\rho_0 = f|\mathrm{GHZ}_N^{+}\rangle\langle\mathrm{GHZ}_N^{+}| + \beta\left( I_D- |\mathrm{GHZ}_N^{+}\rangle\langle\mathrm{GHZ}_N^{+}| \right),
\]

where

\[
\beta=\frac{1-f}{D-1}.
\]

Indeed,

\[
\langle\mathrm{GHZ}_N^{+}|\rho_0|\mathrm{GHZ}_N^{+}\rangle=f.
\]

This is equivalently a mixture with the maximally mixed state,

\[
\rho_0 = V|\mathrm{GHZ}_N^{+}\rangle\langle\mathrm{GHZ}_N^{+}| +(1-V)\frac{I_D}{D},
\]

with visibility

\[
V=\frac{Df-1}{D-1}=f-\beta.
\]

{\color{blue}

These derivations follow from appendix F.1. of the referenced work (are slightly cleaner though). 

}

\subparagraph*{2. Single-qubit dephasing channel}

For a single-qubit off-diagonal matrix element, the Lindblad equation gives

\[
\frac{d}{dt}\rho_{01} = -\left(\gamma+i\omega^{(i)}\right)\rho_{01}.
\]

Hence

\[
\rho_{01}(t) = e^{-\gamma t}e^{-i\omega^{(i)}t}\rho_{01}(0).
\]

Equivalently, the dephasing part is the phase-flip channel

\[
\mathcal E_q(\rho)=q\rho+(1-q)\sigma_z\rho\sigma_z, \qquad q=\frac{1+e^{-\gamma t}}{2},
\]

because it transforms the single-qubit coherence as

\[
|0\rangle\langle1| \longmapsto \bigl[q-(1-q)\bigr]|0\rangle\langle1| = (2q-1)|0\rangle\langle1|.
\]

The global GHZ coherence differs on all $N=nd$ qubits. Independent dephasing therefore gives

\[
|0\rangle^{\otimes N}\langle1|^{\otimes N} \longmapsto (2q-1)^N |0\rangle^{\otimes N}\langle1|^{\otimes N}.
\]

Set

\[
s=(2q-1)^N=(2q-1)^{nd}.
\]

{\color{blue}

Explanation of a standard result in the community. 

}

\subparagraph*{3. Phase encoded by the GHZ state}

The encoding unitary is

\[
U(\boldsymbol{x}) = \exp\left[ -\frac{i}{2} \sum_{i=1}^{d}x_i \sum_{k=1}^{n}\sigma_z^{(i,k)} \right].
\]

Since every node contains $n$ qubits,

\[
U(\boldsymbol{x})|0\rangle^{\otimes N} = e^{-\frac{in}{2}\sum_i x_i}|0\rangle^{\otimes N},
\]

while

\[
U(\boldsymbol{x})|1\rangle^{\otimes N} = e^{+\frac{in}{2}\sum_i x_i}|1\rangle^{\otimes N}.
\]

Thus the relative GHZ phase is

\[
\phi=n\sum_{i=1}^{d}x_i.
\]

Using

\[
\theta_1=\frac{1}{\sqrt d}\sum_{i=1}^{d}x_i,
\]

we obtain

\[
\phi=n\sqrt d\,\theta_1, \qquad \frac{\partial\phi}{\partial\theta_1}=n\sqrt d.
\]

{\color{blue}

Explanation of the accumulated phase during a sole unitary evolution. Identical to my Eq.~\eqref{eq:deph_ghz_phase}.

}

\subparagraph*{4. Parameter-dependent block of the density matrix}

Only the subspace spanned by

\[
|\bar0\rangle=|0\rangle^{\otimes N}, \qquad |\bar1\rangle=|1\rangle^{\otimes N}
\]

contains phase information. In this subspace, after encoding and dephasing,

\[
\rho_{\mathrm{GHZ}}(\theta_1)
=
\frac12
\begin{pmatrix}
f+\beta &
(f-\beta)s\,e^{-i\phi}\\[2mm]
(f-\beta)s\,e^{i\phi} &
f+\beta
\end{pmatrix}.
\]

All remaining $D-2$ eigenvalues equal $\beta$ and are independent of $\theta_1$, so they contribute zero QFI.

The two parameter-relevant eigenvalues are

\[
\lambda_\pm = \frac12\left[ f+\beta\pm(f-\beta)s \right].
\]

For the spectral QFI formula,

\[
\mathcal F_{\theta_1} = \sum_{a,b} \frac{2}{\lambda_a+\lambda_b} \left| \left\langle a\left| \frac{\partial\rho}{\partial\theta_1} \right|b\right\rangle \right|^2,
\]

only the two off-diagonal terms between the $\lambda_+$ and $\lambda_-$ eigenvectors contribute. Their magnitude is

\[
\left| \left\langle +\left| \frac{\partial\rho}{\partial\theta_1} \right|-\right\rangle \right| = \frac{|f-\beta|s}{2} \left|\frac{\partial\phi}{\partial\theta_1}\right|.
\]

Because

\[
\lambda_++\lambda_-=f+\beta,
\]

the QFI becomes

\[
\mathcal F_{\theta_1} = \frac{(f-\beta)^2s^2}{f+\beta} \left(\frac{\partial\phi}{\partial\theta_1}\right)^2.
\]

Substituting $s=(2q-1)^{nd}$ and $\partial_{\theta_1}\phi=n\sqrt d$,

\[
\mathcal F_{\theta_1} = dn^2(2q-1)^{2nd} \frac{(f-\beta)^2}{f+\beta}.
\]

{\color{blue}

The derivations above are correct and I find the choice of notation $f-\beta$ and $f+\beta$ interesting. It makes the equation read better.

}

\subparagraph*{5. Simplification in terms of the GHZ fidelity}

Using $D=2^{nd}$ and $\beta=(1-f)/(D-1)$,

\[
f-\beta = \frac{Df-1}{D-1},
\]

and

\[
f+\beta = \frac{(D-2)f+1}{D-1}.
\]

Therefore,

\[
\frac{(f-\beta)^2}{f+\beta} = \frac{(Df-1)^2} {(D-1)\bigl[(D-2)f+1\bigr]}.
\]

Finally, substitute

\[
D=2^{nd}, \qquad f=Fk^{n-1}.
\]

As a consistency check, when $F=k=q=1$,

\[
\mathcal F_{\theta_1}=dn^2,
\]

which is the ideal GHZ result for the scaled average.

{\color{blue}

The derivations above are correct given the initial assumption about the fidelity function. Though, I find the overall notation ,using for example the $q$ variable, very convoluted, which is unfortunately taken as-is from the peer-reviewed and published work.

}
